\documentclass[12pt,a4paper]{article}
\usepackage[utf8]{inputenc}
\usepackage[english]{babel}
\usepackage{geometry}
\geometry{top=2.5cm, bottom=2.5cm, left=2.5cm, right=2.5cm}
\usepackage{hyperref}
\usepackage{xcolor}
\usepackage{enumitem}
\usepackage{listings}

\title{\textbf{MatricuLAB 2.0} \\ \Large Technical Changelog}
\author{MatricuLAB Development Team}
\date{\today}

\begin{document}

\maketitle

\section*{Overview}
Welcome to MatricuLAB version 2.0. This release introduces fundamental architectural changes to the schedule optimization engine, significant improvements to the UI, and powerful PDF parsing capabilities.

\section{Legacy Features Retained (v1.0)}
Before this update, MatricuLAB served as a manual scheduling utility with the following core functionalities:
\begin{itemize}[label=\textcolor{blue}{\textbullet}]
    \item \textbf{Manual Grid Scheduling:} Visual timetable generation using mock data for manual course placements.
    \item \textbf{Parallel PDF Filtering:} Support for uploading the "Cursos Habilitados" PDF to automatically map and filter courses (Eligible, Mandatory, Elective) against the database.
    \item \textbf{Google Calendar Export:} Full integration to export a processed "Consolidado de Matrícula" or "Horario" PDF directly to the user's Google Calendar account.
\end{itemize}

\section{Major Technical Updates (v2.0)}

\subsection{Algorithmic Evolution: From Brute Force to CSP \& Bitmasks}
The legacy approach (v1.0) relied on computationally heavy brute-force loops and raw string comparisons (\texttt{if (a.day === b.day \&\& a.start < b.end)}). This $O(n^2)$ approach per schedule permutation resulted in severe bottlenecks.
We completely rewrote the "Generate Schedule" core engine to use advanced optimization techniques:
\begin{itemize}[label=\textcolor{blue}{\textbullet}]
    \item \textbf{Constraint Satisfaction Problem (CSP) \& MRV Heuristic:} Instead of evaluating permutations blindly, courses are structured in a hierarchical tree sorted by the \textbf{Minimum Remaining Values (MRV)} heuristic. The algorithm places courses with the fewest available sections (highly constrained) first. If a highly constrained course creates a conflict, the engine prunes the entire branch early, skipping millions of invalid combinations.
    \item \textbf{Hashes \& Bitmasks Collision Detection:} String comparisons were stripped out. The 7-day week is now mapped into integer representations (Hashes) and \textbf{Bitmasks}. Each 30-minute block is a bit in a binary string. Conflict detection is resolved in $O(1)$ time complexity using a simple bitwise AND operation (\texttt{scheduleMask \& sessionMask}).
    \item \textbf{Advanced Mode (N-choose-K Combinatorics):} Users can define a broad pool of potential courses (filtering by eligibles, electives, etc.) and specify a target course count (e.g., "Choose 5 out of 10"). The engine computes $C(n, k)$ subsets and evaluates all nested section permutations across the entire subset space.
    \item \textbf{Scoring System:} Valid schedules are ranked using heuristics based on user preferences (Morning vs Afternoon, minimum gaps, minimum attendance days, lunch break enforcement).
    \item \textbf{UI Tracking:} Added absolute scoring (\texttt{pts}) visualization for transparency over penalty application.
\end{itemize}

\subsection{Web Worker Swarm (Map-Reduce Architecture)}
The schedule combinatorics engine was previously bottlenecked by JavaScript's single-threaded nature.
\begin{itemize}[label=\textcolor{blue}{\textbullet}]
    \item \textbf{Implementation:} The main thread computes the initial combinations domain and chunks the array into $N$ parts based on \texttt{navigator.hardwareConcurrency}.
    \item \textbf{Concurrency:} We instantiate $N$ independent \texttt{Worker} instances (a Swarm). Each worker independently solves its chunk of the CSP problem in parallel.
    \item \textbf{Aggregation:} The main thread acts as a reducer, aggregating \texttt{PROGRESS} events for a unified UI progress bar, and reducing \texttt{COMPLETE} events to merge the Top 200 schedules from all threads into a single globally sorted array.
\end{itemize}

\subsection{'Day Gaps' CSP Penalty}
Users can now penalize "Day Gaps" (e.g., having classes on Mon and Wed, but an empty Tue).
\begin{itemize}[label=\textcolor{blue}{\textbullet}]
    \item \textbf{Metric Tracking:} \texttt{calculateMetricsFromSessions} computes the active span. The mathematical day gap is $span - activeDaysCount$.
    \item \textbf{Heuristic Weighting:} A severe penalty of $-1000$ points per gap day is applied in the worker, ensuring fragmented schedules are immediately discarded.
\end{itemize}

\subsection{Grid UI Redesign}
The manual scheduling grid (\texttt{TimetableGrid.tsx}) received a major facelift. Course block cards now feature a modern, high-contrast aesthetic with rounded corners, drop shadows, and better typographic hierarchy for rapid reading of locations and sections.

\subsection{PDF Parsing Robustness (\texttt{pdfParser.ts})}
The regex parsing logic for UTEC PDFs has been heavily refactored.
\begin{itemize}[label=\textcolor{blue}{\textbullet}]
    \item \textbf{Student Name Extraction:} Corrected the lookahead regex (\texttt{Alumno}). It now successfully breaks on \texttt{Programa}, \texttt{Carrera}, or \texttt{Malla}, preventing greediness.
    \item \textbf{Location Normalization:} Improved \texttt{formatLocation} utility to properly clean \texttt{UTEC-BA} location tags, even with arbitrary spacing.
    \item \textbf{Section Parsing \& Parent Identification:} Resolved recursive section label duplication. Additionally, we implemented intelligent \textbf{parent section identification based on vacancy counts}, correctly grouping complex course structures (like \textit{Química General} in the Excel parser and \textit{Álgebra Lineal} in the PDF) by linking \texttt{Laboratorio} and \texttt{Teoría} sub-groups to their true parent section without errors.
\end{itemize}

\end{document}
